\documentclass[11pt]{article}

\usepackage[margin=1in]{geometry}
\usepackage{amsmath, amssymb}
\usepackage{algorithm}
\usepackage{algpseudocode}
\usepackage{booktabs}
\usepackage{hyperref}
\usepackage{braket}

\title{CSE 4412: Introduction to Quantum Computing - Week 2 Notes}
\author{}
\date{}

\begin{document}
\maketitle
\section*{D-separation and Conditional Independence}
D-separation is a graphical criterion used to determine conditional independence in Bayesian networks. It allows us to identify whether two sets of variables are independent given a third set of variables.
\section*{Bayes ball algorithm}
We can define a set of balls, adn then we try to move then without them colliding ot each other if that move is permitted.

WE cannot write an undirected grphical model such that this is the only conditinoal independece statement, the only conditional independeicies implied here are below. We cannot write a directed graphical model such that these are the only 

defintion, a clique is a completely connected subgraph. A maximal clique is a clique that cannot be extended by including an adjacent vertex. A maximum clique is a clique of the largest possible size in a given graph.
\end{document}
